\section*{Hoofdstuk \ref{ch:b1:elimination} --- \texorpdfstring{$\beta$}{Beta}-eliminatie}

\begin{solution}{exo:b1:elimination:1}
2-Broombutaan: but-1-een, \cip{E}- en \cip{Z}-but-2-een; hoofdproduct
\cip{E}-but-2-een. 2-Broom-2-methyl\-butaan: 2-methylbut-2-een
(hoofdproduct, trigesubstitueerd) en 2-methylbut-1-een. Broomcyclohexaan:
alleen cyclohexeen.
\end{solution}

\begin{solution}{exo:b1:elimination:2}
E1: $v = k[\mathrm{RBr}]$; E2: $v = k[\mathrm{RBr}][\ce{EtO-}]$. Meet de
\omterm{def:b1:rate-laws:initial-rate}{beginsnelheid} van de vorming van het alkeen bij meerdere
ethoxideconcentraties: ze groeit evenredig voor E2, ze verandert niet voor
E1.
\end{solution}

\begin{solution}{exo:b1:elimination:3}
Langs C2--C3 kan het Br op C2 anti staan ten opzichte van elk van beide
waterstofatomen van C3 of ten opzichte van de methylgroep C4: de twee
\omterm{def:b1:stereochemistry-in-depth:isomers}{conformaties} met een H anti ten opzichte van Br kunnen E2 ondergaan (wat
het \cip{E}- en het \cip{Z}-alkeen geeft); de derde kan niet naar C3 toe
elimineren.
\end{solution}

\begin{solution}{exo:b1:elimination:4}
SN2 (primair, hydroxide, lage temperatuur); E2 (tertiair, sterke
omvangrijke base); SN2 (jodide, een goed \omterm{def:b1:electronic-effects:nucleophile}{nucleofiel} en een \omterm{def:b1:predominant-reaction:strong-acid}{zwakke base},
\omterm{def:b1:intermolecular-forces-solvents:solvent-classes}{aprotisch oplosmiddel}); SN1 en E1, waarbij E1 toeneemt bij verwarming.
\end{solution}

\begin{solution}{exo:b1:elimination:5}
C3 draagt één enkel waterstofatoom. In elk \omterm{def:b1:stereochemistry-in-depth:enantiomers}{diastereomeer} legt het
plaatsen van die H anti ten opzichte van het Br van C2 de posities van de
andere groepen vast: de methylgroep van C2 en de ethylgroep van C3 staan
in het ene \omterm{def:b1:stereochemistry-in-depth:enantiomers}{diastereomeer} aan dezelfde kant en in het andere aan
tegengestelde kanten. Het ene geeft \cip{E}-, het andere
\cip{Z}-3-methylpent-2-een: een \omterm{def:b1:nucleophilic-substitution:stereospecific}{stereospecifieke reactie}.
\end{solution}

\begin{solution}{exo:b1:elimination:6}
De OH wordt geprotoneerd, \ce{R-OH2+}; water vertrekt en
er ontstaat het tertiaire \omterm{def:b1:electronic-effects:intermediates}{carbokation} \ce{(CH3)2C+-CH2CH3}; een base (water, \ce{HSO4-})
verwijdert een proton van een naburig koolstofatoom. Hoofdproduct
(Zaitsev): 2-methylbut-2-een.
\end{solution}

\begin{solution}{exo:b1:elimination:7}
De tert-butylgroep vergrendelt de \omterm{def:b1:stereochemistry-in-depth:conformations}{stoel} waarin ze zelf \omterm{def:b1:stereochemistry-in-depth:conformations}{equatoriaal}
staat. In het \textit{trans}-isomeer staat het broom dan ook \omterm{def:b1:stereochemistry-in-depth:conformations}{equatoriaal},
nooit anti ten opzichte van een C--H: E2 moet wachten op de zeer
ongunstige omgeklapte \omterm{def:b1:stereochemistry-in-depth:conformations}{stoel}. In het \textit{cis}-isomeer staat het broom
\omterm{def:b1:stereochemistry-in-depth:conformations}{axiaal}, met twee trans-diaxiale waterstofatomen: snelle E2.
\end{solution}

\begin{solution}{exo:b1:elimination:8}
Ethoxide, klein, verwijdert het waterstofatoom dat het meest
gesubstitueerde alkeen geeft (Zaitsev). tert-Butoxide, omvangrijk,
bereikt de vrijliggende waterstofatomen van de methylgroep gemakkelijker
dan het waterstofatoom van de inwendige \ce{CH2}: het minst
gesubstitueerde alkeen.
\end{solution}

\begin{solution}{exo:b1:elimination:9}
Broommethaan heeft geen $\beta$-koolstofatoom. In
1-broom-2,2-dimethylpropaan is het $\beta$-koolstofatoom quaternair en
draagt het geen waterstof.
\end{solution}

\begin{solution}{exo:b1:elimination:10}
Beide producten hebben het \omterm{def:b1:electronic-effects:intermediates}{carbokation} nodig, dat in de trage stap
gevormd wordt met de snelheid $k_1[\mathrm{RBr}]$. Het wordt daarna
verdeeld tussen invangen door ethanol (\omterm{def:b1:rate-laws:order}{snelheidsconstante} $k_S$) en
verlies van een proton (\omterm{def:b1:rate-laws:order}{snelheidsconstante} $k_E$), beide van
pseudo-eerste orde in het oplosmiddel: de fractie alkeen is
$k_E/(k_E + k_S)$, onafhankelijk van de concentratie van het \omterm{def:b1:nucleophilic-substitution:substitution}{substraat}.
\end{solution}

\begin{solution}{exo:b1:elimination:11}
Vertrek van de \omterm{def:b1:stereochemistry-in-depth:isomers}{conformatie} van de \omterm{def:b1:stereochemistry-in-depth:enantiomers}{mesoverbinding} met de twee Br anti (de
twee fenylgroepen dan anti, de twee H anti), en draai C1 over
$120^\circ$ om zijn H anti ten opzichte van het Br van C2 te plaatsen. De
twee fenylgroepen staan dan aan dezelfde kant van de as H--Br: ze eindigen
\textit{cis}. Op C1 gaat Br vóór Ph; op C2 gaat Ph vóór H; Br en de
fenylgroep van C2 staan \textit{trans}:
\cip{E}-1-broom-1,2-difenyletheen.
\end{solution}

\begin{solution}{exo:b1:elimination:12}
Alleen de \omterm{def:b1:stereochemistry-in-depth:conformations}{stoel} met de \omterm{def:b1:electronic-effects:nucleophile}{vertrekkende groep} \omterm{def:b1:stereochemistry-in-depth:conformations}{axiaal} reageert:
$v = k\,x\,[\mathrm{RY}][\mathrm{B}]$, met $x$ de fractie van die \omterm{def:b1:stereochemistry-in-depth:conformations}{stoel}.
Elke \omterm{def:b1:stereochemistry-in-depth:conformations}{axiale} alkylgroep die ze vereist, kost enkele kJ/mol
(\qty{5.7}{kJ/mol} voor een methylgroep), wat $x$ bij kamertemperatuur
per groep ongeveer door tien deelt: met twee of drie zulke groepen daalt
$x$ tot $10^{-2}$ of $10^{-3}$ en is de reactie evenveel trager.
\end{solution}

\begin{solution}{pb:b1:elimination:1}
\textbf{1.} Een \omterm{def:b1:stereochemistry-in-depth:conformations}{stoel} met Cl, isopropyl en methyl allemaal \omterm{def:b1:stereochemistry-in-depth:conformations}{equatoriaal}.
\textbf{2.} Isopropyl en methyl \omterm{def:b1:stereochemistry-in-depth:conformations}{equatoriaal}, Cl \omterm{def:b1:stereochemistry-in-depth:conformations}{axiaal}.
\textbf{3.} \omterm{def:b1:stereochemistry-in-depth:enantiomers}{Diastereomeren}: ze verschillen alleen bij C1.
\textbf{4.} \omterm{def:b1:electronic-effects:nucleophile}{Vertrekkende groep} en $\beta$-waterstofatoom trans-diaxiaal.
\textbf{5.} Alleen in de omgeklapte \omterm{def:b1:stereochemistry-in-depth:conformations}{stoel}, met Cl, isopropyl en methyl
allemaal \omterm{def:b1:stereochemistry-in-depth:conformations}{axiaal}.
\textbf{6.} Alle drie de groepen \omterm{def:b1:stereochemistry-in-depth:conformations}{axiaal}: de methylgroep alleen kost
ongeveer \qty{5.7}{kJ/mol}, de grotere isopropylgroep meer, het
chlooratoom een beetje: samen zo'n \qty{13}{kJ/mol}, dus een fractie van
ongeveer $\exp(-13000/2479) = \num{5e-3}$, de waarde die het probleem
aanneemt.
\textbf{7.} C2 (één H, het draagt de isopropylgroep) en C6 (twee H).
\textbf{8.} Met Cl \omterm{def:b1:stereochemistry-in-depth:conformations}{axiaal} op C1 de \omterm{def:b1:stereochemistry-in-depth:conformations}{axiale} H van C2 (de isopropylgroep
staat \omterm{def:b1:stereochemistry-in-depth:conformations}{equatoriaal}) en de \omterm{def:b1:stereochemistry-in-depth:conformations}{axiale} H van C6: twee.
\textbf{9.} In de triaxiale \omterm{def:b1:stereochemistry-in-depth:conformations}{stoel} neemt de isopropylgroep de axiale
positie van C2 in, zodat zijn H \omterm{def:b1:stereochemistry-in-depth:conformations}{equatoriaal} staat: niet anti. C6 heeft
een \omterm{def:b1:stereochemistry-in-depth:conformations}{axiale} H: anti.
\textbf{10.} Neomenthylchloride twee, menthylchloride één.
\textbf{11.} Langs C1--C6: Cl (\omterm{def:b1:stereochemistry-in-depth:conformations}{axiaal}, omhoog) op C1 tegenover de \omterm{def:b1:stereochemistry-in-depth:conformations}{axiale}
H (omlaag) op C6; de ringbindingen en de \omterm{def:b1:stereochemistry-in-depth:conformations}{equatoriale} H gestaffeld
daartussen.
\textbf{12.} Op een \omterm{def:b1:stereochemistry-in-depth:conformations}{stoel} staat geen enkele C--H-binding \omterm{def:b1:elimination:periplanar}{synperiplanair}
ten opzichte van de C--Cl-binding; dat afdwingen zou een eclipse,
gespannen ring betekenen.
\textbf{13.} Verwijderen van de H van C2:
4-methyl-1-(propaan-2-yl)cyclohex-1-een (trigesubstitueerd); verwijderen
van een H van C6: 3-methyl-6-(propaan-2-yl)cyclohex-1-een
(digesubstitueerd).
\textbf{14.} Het trigesubstitueerde alkeen, het meest gesubstitueerde en
stabielste.
\textbf{15.} Beide alkenen, \qty{78}{\percent} trigesubstitueerd: in
overeenstemming.
\textbf{16.} Alleen het digesubstitueerde alkeen: het tegendeel van de
voorspelling van Zaitsev, opgelegd door het enige beschikbare
anti-waterstofatoom.
\textbf{17.} Beide zijn stereospecifiek (anti-voorwaarde);
neomenthylchloride reageert regioselectief (78/22); menthylchloride geeft
één enkel constitutie-isomeer.
\textbf{18.} Het koolstofatoom dat Cl draagt, is secundair en geflankeerd
door een isopropylgroep: te vol voor SN2, terwijl ethoxide een \omterm{def:b1:predominant-reaction:strong-acid}{sterke
base} is.
\textbf{19.} $v \propto x\,n_{\mathrm{H}}$: fractie reactieve \omterm{def:b1:stereochemistry-in-depth:conformations}{stoel} maal
het aantal anti-waterstofatomen.
\textbf{20.} $0.95 \times 2 = 1.90$.
\textbf{21.} $\num{5.0e-3} \times 1 = \num{5.0e-3}$.
\textbf{22.} \textbf{$k(\text{neomenthyl})/k(\text{menthyl}) = 1.90/\num{5.0e-3}
= 380$}.
\end{solution}
